\nwfilename{nw/term.nw}\nwbegindocs{0}% -*- mode: poly-noweb; noweb-code-mode: sml-mode; -*-% ===> this file was generated automatically by noweave --- better not edit it
%%
%% Term.nw
%% 
%% Made by Alex Nelson <pqnelson@gmail.com>
%% Login   <alex@lisp>
%% 
%% Started on  2026-04-02T10:39:35-0700
%% Last update 2026-04-02T10:39:35-0700
%% 

\section{Syntax for terms}

\begin{node}[Grammar of terms]
Feferman's $FS_{0}$ requires a three-sorted logic. We can implement
that either as a single-sorted logic with unary predicates for the
sorts (i.e., $\textsf{Ind}[-]$, $\textsf{Fun}[-]$, $\textsf{Class}[-]$),
or we can implement three-sorted terms. We will opt for the latter.

Feferman (\S16) gives the following inductive definitions for the syntax of
these terms:
\begin{enumerate}
\item ``Individuals''
\begin{enumerate}
\item Each ``individual variable'' is in $\textit{InTm}$
\item The constant $\overline{0}$ is in $\textit{InTm}$
\item If $t_{1}\in\textit{InTm}$ and $t_{2}\in\textit{InTm}$, then
  their ordered pair $(t_{1},t_{2})\in\textit{InTm}$
\item If $F\in\textit{FnTm}$ and $t\in\textit{InTm}$, then $Ft\in\textit{InTm}$
\end{enumerate}
\item Functions
\begin{enumerate}
\item Each ``function variable'' is in $\textit{FnTm}$
\item The constants $\operatorname{I}$, $\operatorname{D}$,
  $\operatorname{P}_{1}$, and $\operatorname{P}_{2}$ are in $\textit{FnTm}$
\item If $t\in\textit{InTm}$ then $\operatorname{K}(t)\in\textit{FnTm}$
\item If $F\in\textit{FnTm}$ and $G\in\textit{FnTm}$, then
  $\mathcal{P}(F,G)\in\textit{FnTm}$ and
  $\mathcal{C}(F,G)\in\textit{FnTm}$ and $\mathcal{R}(F,G)\in\textit{FnTm}$
\end{enumerate}
\item Classes
\begin{enumerate}
\item Each ``class variable'' is in $\textit{ClTm}$
\item The singleton $\{\overline{0}\}\in\textit{ClTm}$
\item If $F\in\textit{FnTm}$ and if $S\in\textit{ClTm}$, then the
  preimage $F^{-1}S\in\textit{ClTm}$
\item If $S\in\textit{ClTm}$ and $T\in\textit{ClTm}$, then $S\cup T\in\textit{ClTm}$
  and $S\cap T\in\textit{ClTm}$ and $\mathcal{I}(S,T)\in\textit{ClTm}$
\end{enumerate}
\end{enumerate}
Set comprehension works in a restricted sense (Feferman \S19).

\textbf{Caution:} Note that we need to extend this to support definitions. These would
be ``functions'' (in the usual single-sorted sense --- terms
parametrized by finitely many terms), and each sort would have them.
This makes things really complicated if we leave it at this level of
generality: can we have an individual which depends on a class?
Certainly! We can have a function depend on a class (its
\emph{indicator} function). So we have $\textit{Tm}=\textit{InTm}\cup\textit{FnTm}\cup\textit{ClTm}$
and functions which may depend on finitely many elements of
$\textit{Tm}$. 
\end{node}

\begin{node}
I think the fixed grammar would be:
\begin{enumerate}
\item ``Individuals''
\begin{enumerate}
\item Each ``individual variable'' is in $\textit{InTm}$
\item The constant $\overline{0}$ is in $\textit{InTm}$
\item If $t_{1}\in\textit{InTm}$ and $t_{2}\in\textit{InTm}$, then
  their ordered pair $(t_{1},t_{2})\in\textit{InTm}$
\item If $F\in\textit{FnTm}$ and $t\in\textit{InTm}$, then $Ft\in\textit{InTm}$
\item If $t$ is an arity $n$ constructor of an individual and
  $t_{1},\dots,t_{n}\in\textit{Tm}$, then $t(t_{1},\dots,t_{n})\in\textit{InTm}$
\end{enumerate}
\item Functions
\begin{enumerate}
\item Each ``function variable'' is in $\textit{FnTm}$
\item The constants $\operatorname{I}$, $\operatorname{D}$,
  $\operatorname{P}_{1}$, and $\operatorname{P}_{2}$ are in $\textit{FnTm}$
\item If $t\in\textit{InTm}$ then $\operatorname{K}(t)\in\textit{FnTm}$
\item If $F\in\textit{FnTm}$ and $G\in\textit{FnTm}$, then
  $\mathcal{P}(F,G)\in\textit{FnTm}$ and
  $\mathcal{C}(F,G)\in\textit{FnTm}$ and $\mathcal{R}(F,G)\in\textit{FnTm}$
\item If $F$ is an arity $n$ constructor of a function and
  $t_{1},\dots,t_{n}\in\textit{Tm}$, then $F(t_{1},\dots,t_{n})\in\textit{FnTm}$
\end{enumerate}
\item Classes
\begin{enumerate}
\item Each ``class variable'' is in $\textit{ClTm}$
\item The singleton $\{\overline{0}\}\in\textit{ClTm}$
\item If $F\in\textit{FnTm}$ and if $S\in\textit{ClTm}$, then the
  preimage $F^{-1}S\in\textit{ClTm}$
\item If $S\in\textit{ClTm}$ and $T\in\textit{ClTm}$, then $S\cup T\in\textit{ClTm}$
  and $S\cap T\in\textit{ClTm}$ and $\mathcal{I}(S,T)\in\textit{ClTm}$
\item If $C$ is an arity $n$ constructor of a class and
  $t_{1},\dots,t_{n}\in\textit{Tm}$, then $C(t_{1},\dots,t_{n})\in\textit{ClTm}$
\end{enumerate}
\end{enumerate}
This means we \emph{cannot} have the ``separated'' data structures describing
the three sorts. It must be one giant data structure with 22 data constructors.
The {\Tt{}compare\nwendquote} function requires 127 lines of code, alas.
\end{node}

\nwenddocs{}\nwbegincode{1}\sublabel{NWZKcKR-3RJ9Gq-1}\nwmargintag{{\nwtagstyle{}\subpageref{NWZKcKR-3RJ9Gq-1}}}\moddef{Type for \code{}Term.t\edoc{}~{\nwtagstyle{}\subpageref{NWZKcKR-3RJ9Gq-1}}}\endmoddef\nwstartdeflinemarkup\nwenddeflinemarkup
\nwnotused{Type for [[Term.t]]}\nwendcode{}\nwbegindocs{2}\nwdocspar

\nwenddocs{}\nwbegincode{3}\sublabel{NWZKcKR-3GJZXq-1}\nwmargintag{{\nwtagstyle{}\subpageref{NWZKcKR-3GJZXq-1}}}\moddef{Substitution~{\nwtagstyle{}\subpageref{NWZKcKR-3GJZXq-1}}}\endmoddef\nwstartdeflinemarkup\nwenddeflinemarkup
fun subst x t tm =
  if not(is_var x)
  then raise Fail "\nwlinkedidentc{Term.subst}{NWZKcKR-4XuQKM-1}: first argument must be a variable"
  else if not(are_same_sort x t)
  then raise Fail "\nwlinkedidentc{Term.subst}{NWZKcKR-4XuQKM-1}: variable and replacement must be same sort"
  else (case tm of
   (InVar _) => if x = tm then t else tm
 | (InZero) => tm
 | (InPair(t1,t2)) => Ind.mk_pair(subst x t t1, subst x t t2)
 | (InApp(f, t')) => Ind.mk_app(subst x t f, subst x t t')
 | (InFun(f, args)) => Ind.mk_fun(f, map (subst x t) args)
 | (FnVar _) => if x = tm then t else tm
 | (FnConst(ind)) => Fun.mk_const(subst x t ind)
 | (FnJuxt(f1,f2)) => Fun.mk_juxt(subst x t f1, subst x t f2)
 | (FnComp(f1,f2)) => Fun.mk_comp(subst x t f1, subst x t f2)
 | (FnRec(f1,f2)) => Fun.mk_rec(subst x t f1, subst x t f2)
 | (FnFun(f, args)) => Fun.mk_fun(f, map (subst x t) args)
 | (ClVar _) => if x = tm then t else tm
 | (ClSingleton ind) => Class.mk_singleton(subst x t ind)
 | (ClPreimage(f,c)) => Class.mk_preimage(subst x t f, subst x t c)
 | (ClUnion(c1,c2)) => Class.mk_union(subst x t c1, subst x t c2)
 | (ClIntersection(c1,c2)) => Class.mk_intersection(subst x t c1, subst x t c2)
 | (ClInduct(c1,c2)) => Class.mk_ind(subst x t c1, subst x t c2)
 | (ClFun(f,args)) => Class.mk_fun(f, map (subst x t) args)
 | _ => tm);
\nwnotused{Substitution}\nwidentuses{\\{{\nwixident{Term.subst}}{Term.subst}}}\nwindexuse{\nwixident{Term.subst}}{Term.subst}{NWZKcKR-3GJZXq-1}\nwendcode{}\nwbegindocs{4}\nwdocspar

\begin{node}
We need to decide whether we want to keep the implementation
abstracted away from the user, providing constructors (like
{\Tt{}mk{\_}var\nwendquote}, {\Tt{}mk{\_}app\nwendquote}, etc.) and predicates ({\Tt{}is{\_}var\nwendquote}, {\Tt{}is{\_}fun\nwendquote}, etc.)
and destructuring functions; or if we want to make it all transparent. 
My intuition says transparency may lead to difficulties with
sufficiently clever users.

We will also borrow some names from Clojure. For example,
$\mathcal{P}(f,g)$ will be called ``juxt $f$ $g$'', $\mathcal{C}(f,g)$
is ``comp $f$ $g$''.
\end{node}

\nwenddocs{}\nwbegincode{5}\sublabel{NWZKcKR-2HYRk8-1}\nwmargintag{{\nwtagstyle{}\subpageref{NWZKcKR-2HYRk8-1}}}\moddef{Term.sig~{\nwtagstyle{}\subpageref{NWZKcKR-2HYRk8-1}}}\endmoddef\nwstartdeflinemarkup\nwenddeflinemarkup
signature Term = sig
  exception Dest of string;
  type t;
  val is_ind : t -> bool;
  val is_fun : t -> bool;
  val is_class : t -> bool;
  val is_var : t -> bool;
  val are_same_sort : t -> t -> bool;
  val compare : t * t -> order;
  val compare_list : t list -> t list -> order;
  val eq : t -> t -> bool;
  val fv : t -> t list;
  val subst : t -> t -> t -> t;
  structure Ind : sig
              exception AppE;
              exception PairE;
              val mk_var : string -> t;
              val mk_zero : unit -> t;
              val mk_pair : t * t -> t;
              val mk_app : t * t -> t;
              val mk_fun : string * t list -> t;

              val dest_var : t -> string;
              val dest_pair : t -> t * t;
              val dest_app : t -> t * t;
              val dest_fun : t -> string * t list;

              val is_var : t -> bool;
              val is_zero : t -> bool;
              val is_pair : t -> bool;
              val is_app : t -> bool;
              val is_fun : t -> bool;
            end;
  structure Fun : sig
              exception ConstE;
              exception JuxtE;
              exception CompE;
              exception RecE;
              val mk_var : string -> t;
              (* constants *)
              val id : unit -> t;
              val cond : unit -> t;
              val fst : unit -> t;
              val snd : unit -> t;
              (* constructors *)
              val mk_const : t -> t;
              val mk_juxt : t*t -> t;
              val mk_comp : t*t -> t;
              val mk_rec : t*t -> t;
              val mk_fun : string * t list -> t;

              val dest_var : t -> string;
              val dest_const : t -> t;
              val dest_juxt : t -> t*t;
              val dest_comp : t -> t*t;
              val dest_rec : t -> t*t;
              val dest_fun : t -> string * t list;

              val is_var : t -> bool;
              val is_id : t -> bool;
              val is_cond : t -> bool;
              val is_fst : t -> bool;
              val is_snd : t -> bool;
              val is_const : t -> bool;
              val is_juxt : t -> bool;
              val is_comp : t -> bool;
              val is_rec : t -> bool;
              val is_fun : t -> bool;
            end;
  structure Class : sig
              exception SingletonE;
              exception PreimageE;
              exception UnionE;
              exception IntersectionE;
              exception IndE;
              val mk_var : string -> t;
              val mk_singleton : t -> t;
              val mk_preimage : t * t -> t;
              val mk_union : t * t -> t;
              val mk_intersection : t * t -> t;
              val mk_ind : t * t -> t;
              val mk_fun : string * t list -> t;
              
              val dest_var : t -> string;
              val dest_singleton : t -> t;
              val dest_preimage : t -> t * t;
              val dest_union : t -> t * t;
              val dest_intersection : t -> t * t;
              val dest_ind : t -> t * t;
              val dest_fun : t -> string * t list;

              val is_var : t -> bool;
              val is_singleton : t -> bool;
              val is_preimage : t -> bool;
              val is_union : t -> bool;
              val is_intersection : t -> bool;
              val is_ind : t -> bool;
              val is_fun : t -> bool;
            end;
end;
\nwnotused{Term.sig}\nwendcode{}\nwbegindocs{6}\nwdocspar

\begin{node}
The implementation is a little subtle because we have two
simultaneously defined datatypes: one for individuals, and another for
functions. 
\end{node}

\nwenddocs{}\nwbegincode{7}\sublabel{NWZKcKR-4LsRWC-1}\nwmargintag{{\nwtagstyle{}\subpageref{NWZKcKR-4LsRWC-1}}}\moddef{Term.sml~{\nwtagstyle{}\subpageref{NWZKcKR-4LsRWC-1}}}\endmoddef\nwstartdeflinemarkup\nwenddeflinemarkup
\LA{}Term structure~{\nwtagstyle{}\subpageref{NWZKcKR-P3KAS-1}}\RA{}
\nwnotused{Term.sml}\nwendcode{}\nwbegindocs{8}\nwdocspar

\begin{node}
We have the basic exception for when the user attempts to destructure
a term incorrectly (e.g., they try invoking {\Tt{}\nwlinkedidentq{Term.Ind.dest{\_}var}{NWZKcKR-32OGke-1}\nwendquote} on a
function term, or something). We also have the ``unifying'' encoding
of the {\Tt{}\nwlinkedidentq{Term.t}{NWZKcKR-P3KAS-1}\nwendquote} type.

There are a handful of useful functions, which deserve to be discussed
piecemeal, as well as the substructures for each of the three sorts
appearing in the logic: individuals, functions, and classes.
\end{node}

\nwenddocs{}\nwbegincode{9}\sublabel{NWZKcKR-P3KAS-1}\nwmargintag{{\nwtagstyle{}\subpageref{NWZKcKR-P3KAS-1}}}\moddef{Term structure~{\nwtagstyle{}\subpageref{NWZKcKR-P3KAS-1}}}\endmoddef\nwstartdeflinemarkup\nwusesondefline{\\{NWZKcKR-4LsRWC-1}}\nwenddeflinemarkup
structure Term :> Term = struct
exception Dest of string;

datatype t = InVar of string
           | InZero
           | InPair of t * t
           | InApp of t * t
           | InFun of string * t list
           | FnVar of string
           | FnId
           | FnCond
           | FnFst
           | FnSnd
           | FnConst of t
           | FnJuxt of t * t
           | FnComp of t * t
           | FnRec of t * t
           | FnFun of string * t list
           | ClVar of string
           | ClSingleton of t
           | ClPreimage of t * t
           | ClUnion of t * t
           | ClIntersection of t * t
           | ClInduct of t * t
           | ClFun of string * t list;

\LA{}Term predicates~{\nwtagstyle{}\subpageref{NWZKcKR-4fbveH-1}}\RA{}

\LA{}Comparing terms~{\nwtagstyle{}\subpageref{NWZKcKR-3UxZqA-1}}\RA{}

\LA{}Test if two terms are of the same sort~{\nwtagstyle{}\subpageref{nw@notdef}}\RA{}

\LA{}Free variables in a term~{\nwtagstyle{}\subpageref{NWZKcKR-2nRtky-1}}\RA{}

\LA{}Individual substructure~{\nwtagstyle{}\subpageref{NWZKcKR-2qDNlS-1}}\RA{};

\LA{}Function substructure~{\nwtagstyle{}\subpageref{NWZKcKR-goJMQ-1}}\RA{};

\LA{}Class substructure~{\nwtagstyle{}\subpageref{NWZKcKR-3kqvPP-1}}\RA{};

\LA{}Substituting a term for a variable in a term~{\nwtagstyle{}\subpageref{NWZKcKR-4XuQKM-1}}\RA{}
end
\nwindexdefn{\nwixident{Term.Dest}}{Term.Dest}{NWZKcKR-P3KAS-1}\nwindexdefn{\nwixident{Term.t}}{Term.t}{NWZKcKR-P3KAS-1}\eatline
\nwused{\\{NWZKcKR-4LsRWC-1}}\nwidentdefs{\\{{\nwixident{Term.Dest}}{Term.Dest}}\\{{\nwixident{Term.t}}{Term.t}}}\nwendcode{}\nwbegindocs{10}\nwdocspar
\begin{node}
Substitution recursively works its way down through subterms until it
finds a matching variable, in which case it replaces the variable with
the term {\Tt{}t\nwendquote}. However, before we try doing this, we need to check
that the variable has the same sort as its replacement.
\end{node}

\nwenddocs{}\nwbegincode{11}\sublabel{NWZKcKR-4XuQKM-1}\nwmargintag{{\nwtagstyle{}\subpageref{NWZKcKR-4XuQKM-1}}}\moddef{Substituting a term for a variable in a term~{\nwtagstyle{}\subpageref{NWZKcKR-4XuQKM-1}}}\endmoddef\nwstartdeflinemarkup\nwusesondefline{\\{NWZKcKR-P3KAS-1}}\nwenddeflinemarkup
fun subst x t tm =
  if not(is_var x)
  then raise Fail "\nwlinkedidentc{Term.subst}{NWZKcKR-4XuQKM-1}: first argument must be a variable"
  else if not(are_same_sort x t)
  then raise Fail "\nwlinkedidentc{Term.subst}{NWZKcKR-4XuQKM-1}: variable and replacement must be same sort"
  else (case tm of
   (InVar _) => if x = tm then t else tm
 | (InZero) => tm
 | (InPair(t1,t2)) => Ind.mk_pair(subst x t t1, subst x t t2)
 | (InApp(f, t')) => Ind.mk_app(subst x t f, subst x t t')
 | (InFun(f, args)) => Ind.mk_fun(f, map (subst x t) args)
 | (FnVar _) => if x = tm then t else tm
 | (FnConst(ind)) => Fun.mk_const(subst x t ind)
 | (FnJuxt(f1,f2)) => Fun.mk_juxt(subst x t f1, subst x t f2)
 | (FnComp(f1,f2)) => Fun.mk_comp(subst x t f1, subst x t f2)
 | (FnRec(f1,f2)) => Fun.mk_rec(subst x t f1, subst x t f2)
 | (FnFun(f, args)) => Fun.mk_fun(f, map (subst x t) args)
 | (ClVar _) => if x = tm then t else tm
 | (ClSingleton ind) => Class.mk_singleton(subst x t ind)
 | (ClPreimage(f,c)) => Class.mk_preimage(subst x t f, subst x t c)
 | (ClUnion(c1,c2)) => Class.mk_union(subst x t c1, subst x t c2)
 | (ClIntersection(c1,c2)) => Class.mk_intersection(subst x t c1, subst x t c2)
 | (ClInduct(c1,c2)) => Class.mk_ind(subst x t c1, subst x t c2)
 | (ClFun(f,args)) => Class.mk_fun(f, map (subst x t) args)
 | _ => tm);
\nwindexdefn{\nwixident{Term.subst}}{Term.subst}{NWZKcKR-4XuQKM-1}\eatline
\nwused{\\{NWZKcKR-P3KAS-1}}\nwidentdefs{\\{{\nwixident{Term.subst}}{Term.subst}}}\nwendcode{}\nwbegindocs{12}\nwdocspar
\begin{node}
Determining free variables is done recursively:
\begin{itemize}
\item $FV(x)=\{x\}$
\item $FV(0)=\emptyset$
\item $FV((t_{1},t_{2}))=FV(t_{1})\cup FV(t_{2})$
\item function application: $FV(ft)=FV(f)\cup FV(t)$
\item $FV(\operatorname{I})=FV(\operatorname{D})=FV(\operatorname{P}_{1})=FV(\operatorname{P}_{2})=\emptyset$
\item $FV(\operatorname{K}(t))=FV(t)$
\item $FV(\mathcal{P}(F,G))=FV(F)\cup FV(G)$
\item $FV(\mathcal{C}(F,G))=FV(F)\cup FV(G)$
\item $FV(\mathcal{R}(F,G))=FV(F)\cup FV(G)$
\item $FV(\{0\})=\emptyset$
\item $FV(F^{-1}S)=FV(F)\cup FV(S)$
\item $FV(S\cup T)=FV(S)\cup FV(T)$
\item $FV(S\cap T)=FV(S)\cup FV(T)$
\item $FV(\mathcal{I}(S, T))=FV(S)\cup FV(T)$
\end{itemize}
\end{node}

\nwenddocs{}\nwbegincode{13}\sublabel{NWZKcKR-2nRtky-1}\nwmargintag{{\nwtagstyle{}\subpageref{NWZKcKR-2nRtky-1}}}\moddef{Free variables in a term~{\nwtagstyle{}\subpageref{NWZKcKR-2nRtky-1}}}\endmoddef\nwstartdeflinemarkup\nwusesondefline{\\{NWZKcKR-P3KAS-1}}\nwenddeflinemarkup
fun fv (x as (InVar _)) = [x]
  | fv (InZero) = []
  | fv (InPair(x,y)) = (fv x) @ (fv y)
  | fv (InApp(f,x)) = (fv f) @ (fv x)
  | fv (InFun(_, args)) = List.foldr (fn (arg, fvs) => (fv arg) @ fvs) 
                                     []
                                     args
  | fv (f as (FnVar _)) = [f]
  | fv (FnId) = []
  | fv (FnCond) = []
  | fv (FnFst) = []
  | fv (FnSnd) = []
  | fv (FnConst ind) = fv(ind)
  | fv (FnJuxt(f1,f2)) = (fv f1) @ (fv f2)
  | fv (FnComp(f1,f2)) = (fv f1) @ (fv f2)
  | fv (FnRec(f1,f2)) = (fv f1) @ (fv f2)
  | fv (FnFun(_, args)) = List.foldr (fn (arg, fvs) => (fv arg) @ fvs) 
                                     []
                                     args
  | fv (x as (ClVar _)) = [x]
  | fv (ClSingleton(ind)) = fv ind
  | fv (ClPreimage(f,c)) = (fv f) @ (fv c)
  | fv (ClUnion(c1, c2)) = (fv c1) @ (fv c2)
  | fv (ClIntersection(c1, c2)) = (fv c1) @ (fv c2)
  | fv (ClInduct(c1, c2)) = (fv c1) @ (fv c2)
  | fv (ClFun(_, args)) = List.foldr (fn (arg, fvs) => (fv arg) @ fvs) 
                                     []
                                     args;
\nwindexdefn{\nwixident{Term.fv}}{Term.fv}{NWZKcKR-2nRtky-1}\eatline
\nwused{\\{NWZKcKR-P3KAS-1}}\nwidentdefs{\\{{\nwixident{Term.fv}}{Term.fv}}}\nwendcode{}\nwbegindocs{14}\nwdocspar
\nwenddocs{}\nwbegincode{15}\sublabel{NWZKcKR-4fbveH-1}\nwmargintag{{\nwtagstyle{}\subpageref{NWZKcKR-4fbveH-1}}}\moddef{Term predicates~{\nwtagstyle{}\subpageref{NWZKcKR-4fbveH-1}}}\endmoddef\nwstartdeflinemarkup\nwusesondefline{\\{NWZKcKR-P3KAS-1}}\nwenddeflinemarkup
fun is_ind (InVar _) = true
  | is_ind (InZero) = true
  | is_ind (InPair _) = true
  | is_ind (InApp _) = true
  | is_ind (InFun _) = true
  | is_ind _ = false;

fun is_fun (FnVar _) = true
  | is_fun (FnId) = true
  | is_fun (FnCond) = true
  | is_fun (FnFst) = true
  | is_fun (FnSnd) = true
  | is_fun (FnConst _) = true
  | is_fun (FnJuxt _) = true
  | is_fun (FnComp _) = true
  | is_fun (FnRec _) = true
  | is_fun (FnFun _) = true
  | is_fun _ = false;

fun is_class (ClVar _) = true
  | is_class (ClSingleton _) = true
  | is_class (ClPreimage _) = true
  | is_class (ClUnion _) = true
  | is_class (ClIntersection _) = true
  | is_class (ClInduct _) = true
  | is_class (ClFun _) = true
  | is_class _ = false;

fun are_same_sort lhs rhs =
  if is_ind lhs then is_ind rhs
  else if is_fun lhs then is_fun rhs
  else if is_class lhs then is_class rhs
  else false;

fun is_var (InVar _) = true
  | is_var (FnVar _) = true
  | is_var (ClVar _) = true
  | is_var _ = false;
\nwindexdefn{\nwixident{Term.is{\_}ind}}{Term.is:unind}{NWZKcKR-4fbveH-1}\nwindexdefn{\nwixident{Term.is{\_}fun}}{Term.is:unfun}{NWZKcKR-4fbveH-1}\nwindexdefn{\nwixident{Term.is{\_}class}}{Term.is:unclass}{NWZKcKR-4fbveH-1}\nwindexdefn{\nwixident{Term.is{\_}var}}{Term.is:unvar}{NWZKcKR-4fbveH-1}\nwindexdefn{\nwixident{Term.are{\_}same{\_}sort}}{Term.are:unsame:unsort}{NWZKcKR-4fbveH-1}\eatline
\nwused{\\{NWZKcKR-P3KAS-1}}\nwidentdefs{\\{{\nwixident{Term.are{\_}same{\_}sort}}{Term.are:unsame:unsort}}\\{{\nwixident{Term.is{\_}class}}{Term.is:unclass}}\\{{\nwixident{Term.is{\_}fun}}{Term.is:unfun}}\\{{\nwixident{Term.is{\_}ind}}{Term.is:unind}}\\{{\nwixident{Term.is{\_}var}}{Term.is:unvar}}}\nwendcode{}\nwbegindocs{16}\nwdocspar
\begin{node}
We also have the basic constructors, predicates, and destructuring
functions for individuals.
\end{node}

\nwenddocs{}\nwbegincode{17}\sublabel{NWZKcKR-2qDNlS-1}\nwmargintag{{\nwtagstyle{}\subpageref{NWZKcKR-2qDNlS-1}}}\moddef{Individual substructure~{\nwtagstyle{}\subpageref{NWZKcKR-2qDNlS-1}}}\endmoddef\nwstartdeflinemarkup\nwusesondefline{\\{NWZKcKR-P3KAS-1}}\nwenddeflinemarkup
structure Ind = struct
exception AppE;
exception PairE;

\LA{}Constructors for individuals~{\nwtagstyle{}\subpageref{NWZKcKR-3ejyss-1}}\RA{}

\LA{}Destructuring functions for individuals~{\nwtagstyle{}\subpageref{NWZKcKR-32OGke-1}}\RA{}

\LA{}Predicates for individuals~{\nwtagstyle{}\subpageref{NWZKcKR-2NqhUj-1}}\RA{}
end
\nwindexdefn{\nwixident{Term.Ind.AppE}}{Term.Ind.AppE}{NWZKcKR-2qDNlS-1}\nwindexdefn{\nwixident{Term.Ind.PairE}}{Term.Ind.PairE}{NWZKcKR-2qDNlS-1}\eatline
\nwused{\\{NWZKcKR-P3KAS-1}}\nwidentdefs{\\{{\nwixident{Term.Ind.AppE}}{Term.Ind.AppE}}\\{{\nwixident{Term.Ind.PairE}}{Term.Ind.PairE}}}\nwendcode{}\nwbegindocs{18}\nwdocspar
\begin{node}
Although it has become fashionable to Curry functions everywhere in
functional programming, there is an advantage to having constructors
expect the result of destructuring functions.
\end{node}
\nwenddocs{}\nwbegincode{19}\sublabel{NWZKcKR-3ejyss-1}\nwmargintag{{\nwtagstyle{}\subpageref{NWZKcKR-3ejyss-1}}}\moddef{Constructors for individuals~{\nwtagstyle{}\subpageref{NWZKcKR-3ejyss-1}}}\endmoddef\nwstartdeflinemarkup\nwusesondefline{\\{NWZKcKR-2qDNlS-1}}\nwenddeflinemarkup
(* mk_var : string -> \nwlinkedidentc{Term.t}{NWZKcKR-P3KAS-1} *)
fun mk_var x = (InVar x);

(* mk_zero : unit -> \nwlinkedidentc{Term.t}{NWZKcKR-P3KAS-1} *)
fun mk_zero () = InZero;

(* mk_pair : \nwlinkedidentc{Term.t}{NWZKcKR-P3KAS-1} * \nwlinkedidentc{Term.t}{NWZKcKR-P3KAS-1} -> \nwlinkedidentc{Term.t}{NWZKcKR-P3KAS-1} *)
fun mk_pair (t1, t2) =
  if not(is_ind t1) then raise PairE
  else if not(is_ind t2) then raise PairE
  else InPair(t1,t2);

(* mk_app : \nwlinkedidentc{Term.t}{NWZKcKR-P3KAS-1} * \nwlinkedidentc{Term.t}{NWZKcKR-P3KAS-1} -> \nwlinkedidentc{Term.t}{NWZKcKR-P3KAS-1} *)
fun mk_app(f, t) =
  if not(is_fun f) then raise AppE
  else if not(is_ind t) then raise AppE
  else InApp(f, t);

fun mk_fun(f, args) = InFun(f, args);
\nwindexdefn{\nwixident{Term.Ind.mk{\_}var}}{Term.Ind.mk:unvar}{NWZKcKR-3ejyss-1}\nwindexdefn{\nwixident{Term.Ind.mk{\_}zero}}{Term.Ind.mk:unzero}{NWZKcKR-3ejyss-1}\nwindexdefn{\nwixident{Term.Ind.mk{\_}pair}}{Term.Ind.mk:unpair}{NWZKcKR-3ejyss-1}\nwindexdefn{\nwixident{Term.Ind.mk{\_}app}}{Term.Ind.mk:unapp}{NWZKcKR-3ejyss-1}\eatline
\nwused{\\{NWZKcKR-2qDNlS-1}}\nwidentdefs{\\{{\nwixident{Term.Ind.mk{\_}app}}{Term.Ind.mk:unapp}}\\{{\nwixident{Term.Ind.mk{\_}pair}}{Term.Ind.mk:unpair}}\\{{\nwixident{Term.Ind.mk{\_}var}}{Term.Ind.mk:unvar}}\\{{\nwixident{Term.Ind.mk{\_}zero}}{Term.Ind.mk:unzero}}}\nwidentuses{\\{{\nwixident{Term.t}}{Term.t}}}\nwindexuse{\nwixident{Term.t}}{Term.t}{NWZKcKR-3ejyss-1}\nwendcode{}\nwbegindocs{20}\nwdocspar
\begin{node}
If the user mistakenly tries to destructure a term of a different
syntactic class (e.g., applies {\Tt{}\nwlinkedidentq{Term.Ind.dest{\_}var}{NWZKcKR-32OGke-1}\nwendquote} to a pair), then
a {\Tt{}Dest\nwendquote} exception is propagated with the destructuring function's
name and possibly some debugging information.
\end{node}

\nwenddocs{}\nwbegincode{21}\sublabel{NWZKcKR-32OGke-1}\nwmargintag{{\nwtagstyle{}\subpageref{NWZKcKR-32OGke-1}}}\moddef{Destructuring functions for individuals~{\nwtagstyle{}\subpageref{NWZKcKR-32OGke-1}}}\endmoddef\nwstartdeflinemarkup\nwusesondefline{\\{NWZKcKR-2qDNlS-1}}\nwenddeflinemarkup
(* dest_var : \nwlinkedidentc{Term.t}{NWZKcKR-P3KAS-1} -> string *)
fun dest_var (InVar v) = v
  | dest_var _ = raise Dest "Ind.dest_var: expects a Ind-sorted variable";

(* dest_pair : \nwlinkedidentc{Term.t}{NWZKcKR-P3KAS-1} -> \nwlinkedidentc{Term.t}{NWZKcKR-P3KAS-1} * \nwlinkedidentc{Term.t}{NWZKcKR-P3KAS-1} *)
fun dest_pair(InPair args) = args
  | dest_pair _ = raise Dest "Ind.dest_pair: expects a pair"; 

(* dest_app : \nwlinkedidentc{Term.t}{NWZKcKR-P3KAS-1} -> \nwlinkedidentc{Term.t}{NWZKcKR-P3KAS-1} * \nwlinkedidentc{Term.t}{NWZKcKR-P3KAS-1} *)
fun dest_app(InApp args) = args
  | dest_app _ = raise Dest "Ind.dest_app: expects a function application";

fun dest_fun(InFun args) = args
  | dest_fun _ = raise Dest "Ind.dest_fun: expects a function";
\nwindexdefn{\nwixident{Term.Ind.dest{\_}var}}{Term.Ind.dest:unvar}{NWZKcKR-32OGke-1}\nwindexdefn{\nwixident{Term.Ind.dest{\_}pair}}{Term.Ind.dest:unpair}{NWZKcKR-32OGke-1}\nwindexdefn{\nwixident{Term.Ind.dest{\_}app}}{Term.Ind.dest:unapp}{NWZKcKR-32OGke-1}\eatline
\nwused{\\{NWZKcKR-2qDNlS-1}}\nwidentdefs{\\{{\nwixident{Term.Ind.dest{\_}app}}{Term.Ind.dest:unapp}}\\{{\nwixident{Term.Ind.dest{\_}pair}}{Term.Ind.dest:unpair}}\\{{\nwixident{Term.Ind.dest{\_}var}}{Term.Ind.dest:unvar}}}\nwidentuses{\\{{\nwixident{Term.t}}{Term.t}}}\nwindexuse{\nwixident{Term.t}}{Term.t}{NWZKcKR-32OGke-1}\nwendcode{}\nwbegindocs{22}\nwdocspar
\begin{node}
Predicates for individuals are straightforward tests.
\end{node}

\nwenddocs{}\nwbegincode{23}\sublabel{NWZKcKR-2NqhUj-1}\nwmargintag{{\nwtagstyle{}\subpageref{NWZKcKR-2NqhUj-1}}}\moddef{Predicates for individuals~{\nwtagstyle{}\subpageref{NWZKcKR-2NqhUj-1}}}\endmoddef\nwstartdeflinemarkup\nwusesondefline{\\{NWZKcKR-2qDNlS-1}}\nwenddeflinemarkup
(* is_var : \nwlinkedidentc{Term.t}{NWZKcKR-P3KAS-1} -> bool *)
fun is_var (InVar _) = true
  | is_var _ = false;

(* is_zero : \nwlinkedidentc{Term.t}{NWZKcKR-P3KAS-1} -> bool *)
fun is_zero (InZero) = true
  | is_zero _ = false;

(* is_pair : \nwlinkedidentc{Term.t}{NWZKcKR-P3KAS-1} -> bool *)
fun is_pair (InPair _) = true
  | is_pair _ = false;

(* is_app : \nwlinkedidentc{Term.t}{NWZKcKR-P3KAS-1} -> bool *)
fun is_app (InApp _) = true
  | is_app _ = false;

fun is_fun (InFun _) = true
  | is_fun _ = false;
\nwindexdefn{\nwixident{Term.Ind.is{\_}var}}{Term.Ind.is:unvar}{NWZKcKR-2NqhUj-1}\nwindexdefn{\nwixident{Term.Ind.is{\_}zero}}{Term.Ind.is:unzero}{NWZKcKR-2NqhUj-1}\nwindexdefn{\nwixident{Term.Ind.is{\_}pair}}{Term.Ind.is:unpair}{NWZKcKR-2NqhUj-1}\nwindexdefn{\nwixident{Term.Ind.is{\_}app}}{Term.Ind.is:unapp}{NWZKcKR-2NqhUj-1}\eatline
\nwused{\\{NWZKcKR-2qDNlS-1}}\nwidentdefs{\\{{\nwixident{Term.Ind.is{\_}app}}{Term.Ind.is:unapp}}\\{{\nwixident{Term.Ind.is{\_}pair}}{Term.Ind.is:unpair}}\\{{\nwixident{Term.Ind.is{\_}var}}{Term.Ind.is:unvar}}\\{{\nwixident{Term.Ind.is{\_}zero}}{Term.Ind.is:unzero}}}\nwidentuses{\\{{\nwixident{Term.t}}{Term.t}}}\nwindexuse{\nwixident{Term.t}}{Term.t}{NWZKcKR-2NqhUj-1}\nwendcode{}\nwbegindocs{24}\nwdocspar
\nwenddocs{}\nwbegincode{25}\sublabel{NWZKcKR-goJMQ-1}\nwmargintag{{\nwtagstyle{}\subpageref{NWZKcKR-goJMQ-1}}}\moddef{Function substructure~{\nwtagstyle{}\subpageref{NWZKcKR-goJMQ-1}}}\endmoddef\nwstartdeflinemarkup\nwusesondefline{\\{NWZKcKR-P3KAS-1}}\nwenddeflinemarkup
structure Fun = struct
exception ConstE;
exception JuxtE;
exception CompE;
exception RecE;

fun mk_var x = FnVar x;

fun id () = FnId;

fun cond () = FnCond;
fun fst () = FnFst;
fun snd () = FnSnd;
fun mk_const t = if is_ind t then FnConst t
                 else raise ConstE; 
(* constructors *)
fun mk_juxt (f1, f2) =
  if not(is_fun f1) then raise JuxtE
  else if not(is_fun f2) then raise JuxtE
  else FnJuxt(f1, f2);

fun mk_comp (f1, f2) =
  if not(is_fun f1) then raise CompE
  else if not(is_fun f2) then raise CompE
  else FnComp(f1, f2);

fun mk_rec (f1, f2) =
  if not(is_fun f1) then raise RecE
  else if not(is_fun f2) then raise RecE
  else FnRec(f1,f2);

fun mk_fun (f, args) = FnFun(f, args);

fun dest_var (FnVar x) = x
  | dest_var _ = raise Dest "Fun.dest_var: expects a Fun-sorted variable";

fun dest_const (FnConst ind) = ind
  | dest_const _ = raise Dest "Fun.dest_const: expects a constant";

fun dest_juxt (FnJuxt args) = args
  | dest_juxt _ = raise Dest "Fun.dest_juxt";

fun dest_comp (FnComp args) = args
  | dest_comp _ = raise Dest "Fun.dest_comp";

fun dest_rec (FnRec args) = args
  | dest_rec _ = raise Dest "Fun.dest_rec";

fun dest_fun (FnFun args) = args
  | dest_fun _ = raise Dest "Fun.dest_fun";

fun is_var (FnVar _) = true
  | is_var _ = false;

fun is_id (FnId) = true
  | is_id _ = false;

fun is_cond (FnCond) = true
  | is_cond _ = false;

fun is_fst (FnFst) = true
  | is_fst _ = false;

fun is_snd (FnSnd) = true
  | is_snd _ = false;

fun is_const (FnConst _) = true
  | is_const _ = false;

fun is_juxt (FnJuxt _) = true
  | is_juxt _ = false;

fun is_comp (FnComp _) = true
  | is_comp _ = false;

fun is_rec (FnRec _) = true
  | is_rec _ = false;

fun is_fun (FnFun _) = true
  | is_fun _ = false;
end
\nwused{\\{NWZKcKR-P3KAS-1}}\nwendcode{}\nwbegindocs{26}\nwdocspar

\nwenddocs{}\nwbegincode{27}\sublabel{NWZKcKR-3kqvPP-1}\nwmargintag{{\nwtagstyle{}\subpageref{NWZKcKR-3kqvPP-1}}}\moddef{Class substructure~{\nwtagstyle{}\subpageref{NWZKcKR-3kqvPP-1}}}\endmoddef\nwstartdeflinemarkup\nwusesondefline{\\{NWZKcKR-P3KAS-1}}\nwenddeflinemarkup
structure Class = struct
exception SingletonE;
exception PreimageE;
exception UnionE;
exception IntersectionE;
exception IndE;

fun mk_var x = (ClVar x);

fun mk_singleton(ind) = if is_ind ind then ClSingleton ind
                        else raise SingletonE;

fun mk_preimage (f, c) =
  if not(is_fun f) then raise PreimageE
  else if not(is_class c) then raise PreimageE
  else ClPreimage(f,c);

fun mk_union (c1, c2) =
  if not(is_class c1) orelse not(is_class c2)
  then raise UnionE
  else ClUnion(c1,c2);

fun mk_intersection (c1, c2) =
  if not(is_class c1) orelse not(is_class c2)
  then raise IntersectionE
  else ClIntersection(c1,c2);

fun mk_ind (c1, c2) =
  if not(is_class c1) orelse not(is_class c2)
  then raise IndE
  else ClInduct(c1,c2);

fun mk_fun (f, args) = ClFun(f, args);

fun dest_var (ClVar x) = x
  | dest_var _ = raise Dest "Class.dest_var";

fun dest_singleton (ClSingleton s) = s
  | dest_singleton _ = raise Dest "Class.dest_singleton";

fun dest_preimage (ClPreimage args) = args
  | dest_preimage _ = raise Dest "Class.dest_preimage";

fun dest_union (ClUnion args) = args
  | dest_union _ = raise Dest "Class.dest_union";

fun dest_intersection (ClIntersection args) = args
  | dest_intersection _ = raise Dest "Class.dest_intersection";

fun dest_ind (ClInduct args) = args
  | dest_ind _ = raise Dest "Class.dest_ind";

fun dest_fun (ClFun args) = args
  | dest_fun _ = raise Dest "Class.dest_fun";

fun is_var (ClVar _) = true
  | is_var _ = false;

fun is_singleton (ClSingleton _) = true
  | is_singleton _ = false;

fun is_preimage (ClPreimage _) = true
  | is_preimage _ = false;

fun is_union (ClUnion _) = true
  | is_union _ = false;

fun is_intersection (ClIntersection _) = true
  | is_intersection _ = false;

fun is_ind (ClInduct _) = true
  | is_ind _ = false;

fun is_fun (ClFun _) = true
  | is_fun _ = false;
end
\nwused{\\{NWZKcKR-P3KAS-1}}\nwendcode{}\nwbegindocs{28}\nwdocspar

\begin{node}
Comparing two terms delegates a lot of the work to helper
functions. Broadly speaking: individuals are less than both functions
and classes, and functions are less than classes.
\end{node}

\nwenddocs{}\nwbegincode{29}\sublabel{NWZKcKR-3UxZqA-1}\nwmargintag{{\nwtagstyle{}\subpageref{NWZKcKR-3UxZqA-1}}}\moddef{Comparing terms~{\nwtagstyle{}\subpageref{NWZKcKR-3UxZqA-1}}}\endmoddef\nwstartdeflinemarkup\nwusesondefline{\\{NWZKcKR-P3KAS-1}}\nwenddeflinemarkup
fun compare(InVar x, InVar y) = String.compare(x, y)
  | compare(InVar _, _) = LESS
  | compare(InZero, InVar _) = GREATER
  | compare(InZero, InZero) = EQUAL
  | compare(InZero, _) = LESS
  | compare(InPair _, InVar _) = GREATER
  | compare(InPair _, InZero) = GREATER
  | compare(InPair(a,b), InPair(x,y)) =
     (case compare(a, x) of
       EQUAL => compare(b, y)
     | r => r)
  | compare(InPair _, _) = LESS
  | compare(InApp _, InVar _) = GREATER
  | compare(InApp _, InZero) = GREATER
  | compare(InApp _, InPair _) = GREATER
  | compare(InApp(f,x), InApp(g,y)) =
     (case compare(f, g) of
       EQUAL => compare(x, y)
     | r => r)
  | compare(InApp _, _) = LESS
  | compare(InFun _, InVar _) = GREATER
  | compare(InFun _, InZero) = GREATER
  | compare(InFun _, InPair _) = GREATER
  | compare(InFun _, InApp _) = GREATER
  | compare(InFun(f,args1), InFun(g, args2)) =
     (case String.compare(f,g) of
       EQUAL => compare_list args1 args2
     | r => r)
  | compare(InFun _, _) = LESS
  | compare(FnVar x, FnVar y) = String.compare(x,y)
  | compare(FnVar x, rhs) = if is_ind rhs then GREATER else LESS
  | compare(FnId, FnVar _) = GREATER
  | compare(FnId, FnId) = EQUAL
  | compare(FnId, rhs) = if is_ind rhs then GREATER else LESS
  | compare(FnCond, FnVar _) = GREATER
  | compare(FnCond, FnId) = GREATER
  | compare(FnCond, FnCond) = EQUAL
  | compare(FnCond, rhs) = if is_ind rhs then GREATER else LESS
  | compare(FnFst, FnVar _) = GREATER
  | compare(FnFst, FnId) = GREATER
  | compare(FnFst, FnCond) = GREATER
  | compare(FnFst, FnFst) = EQUAL
  | compare(FnFst, rhs) = if is_ind rhs then GREATER else LESS
  | compare(FnSnd, FnVar _) = GREATER
  | compare(FnSnd, FnId) = GREATER
  | compare(FnSnd, FnCond) = GREATER
  | compare(FnSnd, FnFst) = GREATER
  | compare(FnSnd, FnSnd) = EQUAL
  | compare(FnSnd, rhs) = if is_ind rhs then GREATER else LESS
  | compare(FnConst _, FnVar _) = GREATER
  | compare(FnConst _, FnId) = GREATER
  | compare(FnConst _, FnCond) = GREATER
  | compare(FnConst _, FnFst) = GREATER
  | compare(FnConst _, FnSnd) = GREATER
  | compare(FnConst lhs, FnConst rhs) = compare(lhs, rhs)
  | compare(FnConst _, rhs) = if is_ind rhs then GREATER else LESS
  | compare(FnJuxt _, FnVar _) = GREATER
  | compare(FnJuxt _, FnId) = GREATER
  | compare(FnJuxt _, FnCond) = GREATER
  | compare(FnJuxt _, FnFst) = GREATER
  | compare(FnJuxt _, FnSnd) = GREATER
  | compare(FnJuxt(f1,f2), FnJuxt(g1,g2)) =
     (case compare(f1, g1) of
       EQUAL => compare(f2, g2)
     | r => r)
  | compare(FnJuxt _, rhs) = if is_ind rhs then GREATER else LESS
  | compare(FnComp _, FnVar _) = GREATER
  | compare(FnComp _, FnId) = GREATER
  | compare(FnComp _, FnCond) = GREATER
  | compare(FnComp _, FnFst) = GREATER
  | compare(FnComp _, FnSnd) = GREATER
  | compare(FnComp _, FnJuxt _) = GREATER
  | compare(FnComp(f1,f2), FnComp(g1,g2)) =
     (case compare(f1, g1) of
       EQUAL => compare(f2, g2)
     | r => r)
  | compare(FnComp _, rhs) = if is_ind rhs then GREATER else LESS
  | compare(FnRec _, FnVar _) = GREATER
  | compare(FnRec _, FnId) = GREATER
  | compare(FnRec _, FnCond) = GREATER
  | compare(FnRec _, FnFst) = GREATER
  | compare(FnRec _, FnSnd) = GREATER
  | compare(FnRec _, FnJuxt _) = GREATER
  | compare(FnRec _, FnComp _) = GREATER
  | compare(FnRec(f1,f2), FnRec(g1,g2)) =
     (case compare(f1, g1) of
       EQUAL => compare(f2, g2)
     | r => r)
  | compare(FnRec _, rhs) = if is_ind rhs then GREATER else LESS
  | compare(FnFun _, FnVar _) = GREATER
  | compare(FnFun _, FnId) = GREATER
  | compare(FnFun _, FnCond) = GREATER
  | compare(FnFun _, FnFst) = GREATER
  | compare(FnFun _, FnSnd) = GREATER
  | compare(FnFun _, FnJuxt _) = GREATER
  | compare(FnFun _, FnComp _) = GREATER
  | compare(FnFun _, FnRec _) = GREATER
  | compare(FnFun(f1,args1), FnFun(f2, args2)) =
     (case String.compare(f1, f1) of
       EQUAL => compare_list args1 args2
     | r => r)
  | compare(FnFun _, rhs) = if is_ind rhs then GREATER else LESS
  | compare(ClVar x, ClVar y) = String.compare(x, y)
  | compare(ClVar _, rhs) = if is_class rhs then LESS else GREATER
  | compare(ClSingleton _, ClVar _) = GREATER
  | compare(ClSingleton t1, ClSingleton t2) = compare(t1, t2)
  | compare(ClSingleton _, rhs) = if is_class rhs then LESS else GREATER
  | compare(ClPreimage _, ClVar _) = GREATER
  | compare(ClPreimage _, ClSingleton _) = GREATER
  | compare(ClPreimage(f1,c1), ClPreimage(f2,c2)) =
     (case compare(f1, f2) of
       EQUAL => compare(c1, c2)
     | r => r)
  | compare(ClPreimage _, rhs) = if is_class rhs then LESS else GREATER
  | compare(ClUnion _, ClVar _) = GREATER
  | compare(ClUnion _, ClSingleton _) = GREATER
  | compare(ClUnion _, ClPreimage _) = GREATER
  | compare(ClUnion(c1,c2), ClUnion(d1,d2)) =
     (case compare(c1, d1) of
       EQUAL => compare(c2, d2)
     | r => r)
  | compare(ClUnion _, rhs) = if is_class rhs then LESS else GREATER
  | compare(ClIntersection _, ClVar _) = GREATER
  | compare(ClIntersection _, ClSingleton _) = GREATER
  | compare(ClIntersection _, ClPreimage _) = GREATER
  | compare(ClIntersection _, ClUnion _) = GREATER
  | compare(ClIntersection(c1,c2), ClIntersection(d1,d2)) =
     (case compare(c1, d1) of
       EQUAL => compare(c2, d2)
     | r => r)
  | compare(ClIntersection _, rhs) = if is_class rhs then LESS else GREATER
  | compare(ClInduct _, ClVar _) = GREATER
  | compare(ClInduct _, ClSingleton _) = GREATER
  | compare(ClInduct _, ClPreimage _) = GREATER
  | compare(ClInduct _, ClUnion _) = GREATER
  | compare(ClInduct _, ClIntersection _) = GREATER
  | compare(ClInduct(c1,c2), ClInduct(d1,d2)) =
     (case compare(c1, d1) of
       EQUAL => compare(c2, d2)
     | r => r)
  | compare(ClInduct _, rhs) = if is_class rhs then LESS else GREATER
  | compare(ClFun(f1,args1), ClFun(f2,args2)) = 
     (case String.compare(f1, f2) of
       EQUAL => compare_list args1  args2
     | r => r)
  | compare(ClFun _, _) = GREATER
and compare_list [] [] = EQUAL
  | compare_list [] ys = LESS
  | compare_list xs [] = GREATER
  | compare_list (x::xs) (y::ys) =
     (case compare(x, y) of
       EQUAL => compare_list xs ys
     | r => r);

fun eq x y = EQUAL = compare(x,y);
\nwindexdefn{\nwixident{Term.compare}}{Term.compare}{NWZKcKR-3UxZqA-1}\eatline
\nwused{\\{NWZKcKR-P3KAS-1}}\nwidentdefs{\\{{\nwixident{Term.compare}}{Term.compare}}}\nwendcode{}\nwbegindocs{30}\nwdocspar
\nwenddocs{}

\nwixlogsorted{c}{{Class substructure}{NWZKcKR-3kqvPP-1}{\nwixu{NWZKcKR-P3KAS-1}\nwixd{NWZKcKR-3kqvPP-1}}}%
\nwixlogsorted{c}{{Comparing terms}{NWZKcKR-3UxZqA-1}{\nwixu{NWZKcKR-P3KAS-1}\nwixd{NWZKcKR-3UxZqA-1}}}%
\nwixlogsorted{c}{{Constructors for individuals}{NWZKcKR-3ejyss-1}{\nwixu{NWZKcKR-2qDNlS-1}\nwixd{NWZKcKR-3ejyss-1}}}%
\nwixlogsorted{c}{{Destructuring functions for individuals}{NWZKcKR-32OGke-1}{\nwixu{NWZKcKR-2qDNlS-1}\nwixd{NWZKcKR-32OGke-1}}}%
\nwixlogsorted{c}{{Free variables in a term}{NWZKcKR-2nRtky-1}{\nwixu{NWZKcKR-P3KAS-1}\nwixd{NWZKcKR-2nRtky-1}}}%
\nwixlogsorted{c}{{Function substructure}{NWZKcKR-goJMQ-1}{\nwixu{NWZKcKR-P3KAS-1}\nwixd{NWZKcKR-goJMQ-1}}}%
\nwixlogsorted{c}{{Individual substructure}{NWZKcKR-2qDNlS-1}{\nwixu{NWZKcKR-P3KAS-1}\nwixd{NWZKcKR-2qDNlS-1}}}%
\nwixlogsorted{c}{{Predicates for individuals}{NWZKcKR-2NqhUj-1}{\nwixu{NWZKcKR-2qDNlS-1}\nwixd{NWZKcKR-2NqhUj-1}}}%
\nwixlogsorted{c}{{Substituting a term for a variable in a term}{NWZKcKR-4XuQKM-1}{\nwixu{NWZKcKR-P3KAS-1}\nwixd{NWZKcKR-4XuQKM-1}}}%
\nwixlogsorted{c}{{Substitution}{NWZKcKR-3GJZXq-1}{\nwixd{NWZKcKR-3GJZXq-1}}}%
\nwixlogsorted{c}{{Term predicates}{NWZKcKR-4fbveH-1}{\nwixu{NWZKcKR-P3KAS-1}\nwixd{NWZKcKR-4fbveH-1}}}%
\nwixlogsorted{c}{{Term structure}{NWZKcKR-P3KAS-1}{\nwixu{NWZKcKR-4LsRWC-1}\nwixd{NWZKcKR-P3KAS-1}}}%
\nwixlogsorted{c}{{Term.sig}{NWZKcKR-2HYRk8-1}{\nwixd{NWZKcKR-2HYRk8-1}}}%
\nwixlogsorted{c}{{Term.sml}{NWZKcKR-4LsRWC-1}{\nwixd{NWZKcKR-4LsRWC-1}}}%
\nwixlogsorted{c}{{Test if two terms are of the same sort}{nw@notdef}{\nwixu{NWZKcKR-P3KAS-1}}}%
\nwixlogsorted{c}{{Type for \code{}Term.t\edoc{}}{NWZKcKR-3RJ9Gq-1}{\nwixd{NWZKcKR-3RJ9Gq-1}}}%
\nwixlogsorted{i}{{\nwixident{Term.are{\_}same{\_}sort}}{Term.are:unsame:unsort}}%
\nwixlogsorted{i}{{\nwixident{Term.compare}}{Term.compare}}%
\nwixlogsorted{i}{{\nwixident{Term.Dest}}{Term.Dest}}%
\nwixlogsorted{i}{{\nwixident{Term.fv}}{Term.fv}}%
\nwixlogsorted{i}{{\nwixident{Term.Ind.AppE}}{Term.Ind.AppE}}%
\nwixlogsorted{i}{{\nwixident{Term.Ind.dest{\_}app}}{Term.Ind.dest:unapp}}%
\nwixlogsorted{i}{{\nwixident{Term.Ind.dest{\_}pair}}{Term.Ind.dest:unpair}}%
\nwixlogsorted{i}{{\nwixident{Term.Ind.dest{\_}var}}{Term.Ind.dest:unvar}}%
\nwixlogsorted{i}{{\nwixident{Term.Ind.is{\_}app}}{Term.Ind.is:unapp}}%
\nwixlogsorted{i}{{\nwixident{Term.Ind.is{\_}pair}}{Term.Ind.is:unpair}}%
\nwixlogsorted{i}{{\nwixident{Term.Ind.is{\_}var}}{Term.Ind.is:unvar}}%
\nwixlogsorted{i}{{\nwixident{Term.Ind.is{\_}zero}}{Term.Ind.is:unzero}}%
\nwixlogsorted{i}{{\nwixident{Term.Ind.mk{\_}app}}{Term.Ind.mk:unapp}}%
\nwixlogsorted{i}{{\nwixident{Term.Ind.mk{\_}pair}}{Term.Ind.mk:unpair}}%
\nwixlogsorted{i}{{\nwixident{Term.Ind.mk{\_}var}}{Term.Ind.mk:unvar}}%
\nwixlogsorted{i}{{\nwixident{Term.Ind.mk{\_}zero}}{Term.Ind.mk:unzero}}%
\nwixlogsorted{i}{{\nwixident{Term.Ind.PairE}}{Term.Ind.PairE}}%
\nwixlogsorted{i}{{\nwixident{Term.is{\_}class}}{Term.is:unclass}}%
\nwixlogsorted{i}{{\nwixident{Term.is{\_}fun}}{Term.is:unfun}}%
\nwixlogsorted{i}{{\nwixident{Term.is{\_}ind}}{Term.is:unind}}%
\nwixlogsorted{i}{{\nwixident{Term.is{\_}var}}{Term.is:unvar}}%
\nwixlogsorted{i}{{\nwixident{Term.subst}}{Term.subst}}%
\nwixlogsorted{i}{{\nwixident{Term.t}}{Term.t}}%

